\section*{अध्याय \ref{ch:b1:radiometric-methods} --- रेडियोमेट्रिक आयु-निर्धारण की विधियाँ}
% values: calc/ch25.py

\begin{solution}{exo:b1:radiometric-methods:1}
(क)~$D^*/P = 2^n - 1 = 3$ से $n = 2$: दो अर्ध-आयुएँ। (ख)~रूबिडियम-87: $e^{\lambda t} - 1 = e^{\num{1.3972e-11} \times \num{3.10e8}} - 1 = 0.004341$; यूरेनियम-235: $e^{\num{9.8485e-10} \times \num{3.10e8}} - 1 = 0.3570$। यूरेनियम का अनुपात लगभग 80 गुना बड़ा है: मापन की एक दी हुई निरपेक्ष त्रुटि उसका कहीं छोटा भाग है। (ग)~$2^{-n} < 0.01$ जब $n > \log_2 100 = 6.64$ हो: लगभग सात अर्ध-आयुओं के बाद।
\end{solution}

\begin{solution}{exo:b1:radiometric-methods:2}
(क)~कोयले पर रेडियोकार्बन। (ख)~ज़िरकॉन पर यूरेनियम–सीसा, या सैनिडीन पर आर्गन–आर्गन (पोटैशियम–आर्गन), जो एक पोटैशियम-समृद्ध \omterm{def:b1:rock-forming-minerals:mineral}{खनिज} है। (ग)~ज़िरकॉन पर यूरेनियम–सीसा, जो एक लंबे इतिहास से बचता है; पूरी शैलों पर रूबिडियम–स्ट्रॉन्शियम संभव है; पोटैशियम–आर्गन में बाद के ऊष्मन में खोए आर्गन से आयुएँ घटने का जोखिम होगा। (घ)~सीप के कार्बोनेट पर रेडियोकार्बन, विधि की सीमा के निकट (लगभग \num{70000} वर्ष)। (ङ)~सीधे कोई नहीं: उसकी आयु ऊपर और नीचे की आँकी गई ज्वालामुखीय शैलों से परिबद्ध है, या उसके जीवाश्मों को एक आँकी गई खंड से सहसंबंधित करके मिलती है (\cref{ch:b1:stratigraphic-scale})।
\end{solution}

\begin{solution}{exo:b1:radiometric-methods:3}
(क)~$t = -8033 \ln 0.25 = 8033 \times 1.3863 = \num{11136}$ वर्ष BP। (ख)~माध्य-जीवन $5700/\ln 2 = \num{8223}$ वर्ष है: $t = 8223 \times 1.3863 =
\num{11400}$ वर्ष। (ग)~दोनों मानते हैं कि जीवित कार्बन में कार्बन-14 का अनुपात तब भी वही था जो मानक में है; वह नहीं था। कैलेंडर आयु IntCal20 जैसे एक \omterm{def:b1:radiometric-methods:radiocarbon}{अंशांकन वक्र} से पढ़ी जाती है, cal BP वर्षों की एक परास के रूप में।
\end{solution}

\begin{solution}{exo:b1:radiometric-methods:4}
(क)~हर घड़ी तब शुरू होती है जब उसका \omterm{def:b1:rock-forming-minerals:mineral}{खनिज} अपने \omterm{def:b1:radiometric-methods:closure}{अवरोध-तापमान} से नीचे ठंडा होता है: ज़िरकॉन (\qty{750}{\celsius} से ऊपर) पहले बंद होता है, हॉर्नब्लेंड (लगभग \qty{530}{\celsius}) उसके बाद, बायोटाइट (लगभग \qty{280}{\celsius}) अंत में, इसलिए जैसे-जैसे ग्रेनाइट ठंडा हुआ, आयुएँ उसी क्रम में घटती हैं। (ख)~$(530 - 280)/(410 - 398) =
250/12 = \qty{21}{\celsius}$ प्रति दस लाख वर्ष। (ग)~ज़िरकॉन \qty{750}{\celsius} से ऊपर बंद हुआ, इसलिए ग्रेनाइट कम से कम $750 - 530 =
\qty{220}{\celsius}$ \qty{10}{\Myr} में ठंडा हुआ: कम से कम \qty{22}{\celsius} प्रति दस लाख वर्ष, एक निम्न सीमा।
\end{solution}

\begin{solution}{exo:b1:radiometric-methods:5}
(क)~$\bar x = 7.5$, $\bar y = 0.809525$, $S_{xx} = 81.0$, $S_{xy} = 1.13965$; $m = 1.13965/81.0 = 0.0140698$, $b = 0.809525 - 7.5 \times 0.0140698 = 0.70400$। (ख)~$t = \ln(1.0140698)/\num{1.3972e-11} = \qty{1000.0}{\Ma}$; \omterm{def:b1:radiometric-methods:isochron}{आरंभिक अनुपात} $0.7040$। (ग)~अवशेष हैं $+1.6$, $-2.5$, $+0.7$ और $+0.2$ (\num{e-4} की इकाइयों में): $s = \sqrt{\num{9.34e-8}/2} = 0.00022$, $\sigma_m = 0.00022/\sqrt{81}
= \num{2.40e-5}$ और $\sigma_t = \num{2.40e-5}/(\num{1.3972e-11} \times 1.0141) =
\qty{1.7}{\Myr}$: \qty[uncertainty-mode=separate]{1000.0(17)}{\Ma}। (घ)~$0.7040$ मैंटल के मान $0.7025$ के निकट है और पपड़ी के $0.710$ से काफ़ी नीचे: एक स्रोत जो अधिकतर मैंटल में है।
\end{solution}

\begin{solution}{exo:b1:radiometric-methods:6}
(क)~$0.100/39.098 \times \num{1.17e-4} = \qty{2.992e-7}{mol}$ पोटैशियम-40 प्रति ग्राम। (ख)~${}^{40}\mathrm{Ar}^*/{}^{40}\mathrm{K} = \num{1.169e-9}/\num{2.992e-7}
= 0.0039065$; $\lambda/\lambda_e = 5.543/0.581 = 9.5404$; $t = \ln(1 + 9.5404
\times 0.0039065)/\num{5.543e-10} = \ln(1.037269)/\num{5.543e-10} =
\qty{66.01}{\Ma}$। (ग)~विकिरण-जनित आर्गन $0.995 \times 1.169$ नैनोमोल होता, जो \qty{65.69}{\Ma} देता: मापी गई आयु \qty{0.32}{\Myr} अधिक पुरानी होती, वह \num{32000} वर्ष का लगभग दस गुना, जिनके भीतर प्रभाव और क्रिटेशियस के अंत का विलोपन अब समकालिक दिखाए गए हैं। अतिरिक्त आर्गन मुख्यतः युवा शैलों और संपूर्ण-शैल नमूनों के लिए चिंता का विषय है; सैनिडीन जैसे \omterm{def:b1:rock-forming-minerals:mineral}{खनिज} को क्रिस्टलित होते समय आर्गन नहीं लेना चाहिए, और एक आर्गन–आर्गन आयु-स्पेक्ट्रम इसे प्रकट कर देगा।
\end{solution}

\begin{solution}{exo:b1:radiometric-methods:7}
(क)~लगभग \num{2000} cal BP: वक्र वहाँ \num{2070} वर्ष BP से गुज़रता है। (ख)~लगभग \num{2495} से \num{2699} cal BP तक, दो सौ वर्ष: वक्र \num{2500} और \num{2700} cal BP के बीच लगभग सपाट है, इसलिए उस अंतराल की हर कैलेंडर आयु की \omterm{def:b1:radiometric-methods:radiocarbon}{रेडियोकार्बन आयु} लगभग एक ही है, और अधिक सटीक माप परास को संकीर्ण नहीं करेगा। (ग)~\num{2000} cal BP पर \omterm{def:b1:radiometric-methods:radiocarbon}{रेडियोकार्बन आयु} कैलेंडर आयु से \num{70} वर्ष \emph{पुरानी} है; \num{8000} cal BP पर \num{822} वर्ष छोटी (\qty{10}{\percent}); \num{12000} cal BP पर \num{1694} वर्ष छोटी (\qty{14}{\percent})। अंतर आयु के अनुपात में नहीं है: वह चिह्न भी बदलता है। वह अकेले अर्ध-आयु से नहीं आ सकता, जो एक स्थिर \qty{2.4}{\percent} देती।
\end{solution}

\begin{solution}{exo:b1:radiometric-methods:8}
(क)~नीचे की राख-परत (\qty{94.50}{\Ma}) एक उच्चतम आयु देती है, ऊपर का लावा-प्रवाह (\qty{93.80}{\Ma}) एक न्यूनतम: सीमा उनके बीच परिबद्ध है। अंतर्वेधन भी एक न्यूनतम है, पर कमज़ोर, लावा से युवा। डेट्राइटल ज़िरकॉन केवल बलुआ पत्थर के लिए एक उच्चतम है, और राख-परत से कमज़ोर, जो बलुआ पत्थर के नीचे है। (ख)~$f = (12 - 4)/(20 - 4) = 0.5$; $t = 94.50 - 0.5 \times 0.70 =
\qty{94.15}{\Ma}$; $\sigma_t = \sqrt{0.25 \times 0.10^2 + 0.25 \times 0.12^2} =
\qty{0.08}{\Myr}$: \qty[uncertainty-mode=separate]{94.15(8)}{\Ma}। (ग)~उसके कण कम से कम \qty{95.3}{\Ma} पुरानी एक शैल से आए थे, स्रोत-क्षेत्र में कहीं अपरदित; यह दिखाता है कि बलुआ पत्थर \qty{95.3}{\Ma} से युवा है, जिसे नीचे की राख-परत पहले ही अधिक कसकर सिद्ध कर चुकी है।
\end{solution}

\begin{solution}{exo:b1:radiometric-methods:9}
(क)~कण~A: $\ln(1.08401)/\num{1.55125e-10} = \qty{520.0}{\Ma}$ और $\ln(1.6688)/\num{9.8485e-10} = \qty{520.0}{\Ma}$। कण~B: \qty{368.4}{\Ma} और \qty{389.9}{\Ma}। (ख)~कण~A संगत है। (ग)~$0.05881/0.08401 = 0.7000$ और $0.4682/0.6688 = 0.7001$। आयु $\ln(1 + R)/\lambda$ है, $R$ के अनुपात में नहीं: छोटे यूरेनियम-238 अनुपात के लिए $\ln(1+R) \approx R$ और आयु लगभग \qty{30}{\percent} गिरती है (\qty{368}{\Ma} तक); बड़े यूरेनियम-235 अनुपात के लिए लघुगणक कम गिरता है, और आयु केवल \qty{390}{\Ma} तक। दोनों बहुत युवा हैं, और वे असहमत हैं: कण असंगत है।
\end{solution}

\begin{solution}{exo:b1:radiometric-methods:10}
(क)~$u = \lambda t$ के साथ, $D^*/P = e^u - 1 = u + u^2/2 + \dots$, इसलिए $(D^*/P)/\lambda
= t\,(1 + u/2 + \dots)$: सन्निकटन आयु को लगभग $u/2$ के निकट सापेक्ष त्रुटि से अधिक आँकता है। (ख)~\qty{310}{\Ma} पर रूबिडियम-87: $u = 0.00433$, त्रुटि \qty{0.22}{\percent} (\qty{0.7}{\Myr}); \qty{296}{\Ma} पर पोटैशियम-40: $u =
0.164$, त्रुटि \qty{8.7}{\percent} (\qty{26}{\Myr}; $u/2$ \qty{8.2}{\percent} देता है); \qty{4.4}{\Ga} पर यूरेनियम-238: $u = 0.683$, त्रुटि \qty{43}{\percent} --- सन्निकटन, और उसका त्रुटि-आकलन ($u/2 = \qty{34}{\percent}$) भी, विफल होते हैं। (ग)~$R = 0.004341$, $\sigma_R = \num{4.34e-6}$; \cref{prop:b1:radiometric-methods:precision} से, $\sigma_t =
\num{4.34e-6}/(\num{1.3972e-11} \times 1.0043) = \qty{0.31}{\Myr}$।
\end{solution}

\begin{solution}{exo:b1:radiometric-methods:11}
(क)~मान लीजिए ${}^{40}\mathrm{K}$ आज का पोटैशियम-40 है। वर्तमान से पहले समय $t_2$ पर ${}^{40}\mathrm{K}\,e^{\lambda t_2}$ परमाणु थे, और $t_1$ और $t_2$ के बीच बना आर्गन, उस अंतराल पर लागू \cref{prop:b1:radiometric-methods:kar} से, $(\lambda_e/\lambda)\,{}^{40}\mathrm{K}\,e^{\lambda
t_2}\bigl(e^{\lambda (t_1 - t_2)} - 1\bigr) = (\lambda_e/\lambda)\,
{}^{40}\mathrm{K}\,\bigl(e^{\lambda t_1} - e^{\lambda t_2}\bigr)$ था। उसका अंश $1 - f$ बचा। $t_2$ से, बंद \omterm{def:b1:rock-forming-minerals:mineral}{खनिज} ने $(\lambda_e/\lambda)\,{}^{40}\mathrm{K}\,(e^{\lambda t_2} - 1)$ और बनाया है। जोड़कर और $\lambda/(\lambda_e\,{}^{40}\mathrm{K})$ से गुणा करने पर व्यंजक मिलता है। (ख)~स्पष्ट आयु $\ln(1 + \text{व्यंजक})/\lambda$ है: $f = 0$ के लिए \qty{400.0}{\Myr}, $0.3$ के लिए \qty{342.3}{\Myr}, $0.6$ के लिए \qty{282.7}{\Myr} और $1$ के लिए \qty{200.0}{\Myr}। (ग)~\qty{200}{\Myr} और \qty{400}{\Myr} के बीच की हर आयु किसी आंशिक हानि से बन सकती है, और एक अकेली संख्या नहीं बता सकती कि वह उस समय की एक वास्तविक घटना दर्ज करती है या मिश्रण; एक आयु-स्पेक्ट्रम, या उसी शैल की दूसरी घड़ियाँ, आवश्यक हैं।
\end{solution}

\begin{solution}{exo:b1:radiometric-methods:12}
(क)~आयु प्रयुक्त माध्य-जीवन के अनुपात में है: $10306 \times 8223/8033 =
\num{10550}$ वर्ष। (ख)~$10550 - 10306 = 244$ वर्ष, \num{1694} के लगभग सातवें भाग। (ग)~यदि आरंभिक सक्रियता $k A_0$ थी, तो सही आयु $t = \tau
\ln(k A_0/A) = \num{10550} + \tau \ln k$ है; $t = \num{12000}$ के साथ, $\ln k =
1450/8223 = 0.176$ और $k = 1.19$: मानक से लगभग \qty{19}{\percent} अधिक कार्बन-14।
\end{solution}

\begin{solution}{pb:b1:radiometric-methods:1}
\begin{enumerate}
  \item बिंदु $x$ के साथ लगातार ऊपर उठते हैं; रेखा से उनके विचलन चौथे दशमलव में अधिकतम दो इकाइयाँ हैं, जबकि परास $0.04$ है, $y$ में: वे थोड़े बिखराव के साथ एक सीधी रेखा पर हैं।
  \item $\bar x = 4.90333$, $\bar y = 0.729288$, $S_{xx} = 62.1253$, $S_{xy} =
        0.269335$; $m = 0.269335/62.1253 = 0.0043354$।
  \item $b = 0.729288 - 0.0043354 \times 4.90333 = 0.70803$: ग्रेनाइट का $^{87}\mathrm{Sr}/^{86}\mathrm{Sr}$ जब वह क्रिस्टलित हुआ, उसके सभी भागों द्वारा साझा।
  \item $t = \ln(1.0043354)/\num{1.3972e-11} = \qty{309.6}{\Ma}$।
  \item अवशेष हैं $+9.1$, $-17.2$, $+6.9$, $-6.4$, $+15.1$ और $-7.5$ (\num{e-5} की इकाइयों में); $s = 0.000137$, $\sigma_m = 0.000137/\sqrt{62.1253} =
        \num{1.74e-5}$, $\sigma_t = \num{1.74e-5}/(\num{1.3972e-11} \times 1.0043)
        = \qty{1.2}{\Myr}$: \qty[uncertainty-mode=separate]{309.6(12)}{\Ma}, और $\sigma_b = 0.00010$।
  \item $0.7080$ पपड़ी के $0.710$ के निकट है, मैंटल के $0.7025$ के मुकाबले कहीं अधिक। मिश्रण के लिए, $(0.70803 - 0.7025)/(0.710 - 0.7025) = 0.74$: स्ट्रॉन्शियम का लगभग तीन-चौथाई पिघली पपड़ी से आया।
  \item $0.0740/39.098 \times \num{1.17e-4} = \qty{2.214e-7}{mol}$ प्रति ग्राम, $\num{6.022e23}$ से गुणा करने पर: \num{1.33e17} परमाणु।
  \item $^{40}\mathrm{Ar}^*/^{40}\mathrm{K} = \num{4.139e-9}/\num{2.2144e-7} =
        0.018691$; $1 + 9.5404 \times 0.018691 = 1.17832$; $t =
        \ln(1.17832)/\num{5.543e-10} = \qty{296.0}{\Ma}$।
  \item $^{40}\mathrm{K} = \qty{3.292e-8}{mol.g^{-1}}$, अनुपात $0.019239$, $\ln(1.18355)/\num{5.543e-10} = \qty{304.0}{\Ma}$।
  \item $\ln(1.018691)/\num{5.543e-10} = \qty{33.4}{\Ma}$, नौ गुना बहुत युवा, क्योंकि लगभग हर दस क्षयों में केवल एक आर्गन देता है। कैल्शियम-40 कैल्शियम का सामान्य समस्थानिक है, शैलों में प्रचुर: क्षय से बना थोड़ा, खनिजों के आरंभ से मौजूद कैल्शियम से अलग नहीं किया जा सकता।
  \item अतिरिक्त आर्गन आयुओं को लगभग दस लाख वर्षों की कोटि में खिसकाता है, जो युवा शैलों के लिए मायने रखता है; यहाँ वह आयु के \qty{1}{\percent} से बहुत कम होगा। और \omterm{def:b1:rock-forming-minerals:groups}{अभ्रक} और \omterm{def:b1:rock-forming-minerals:groups}{एम्फ़ीबोल} के क्रिस्टल बनते समय लगभग कोई आर्गन नहीं लेते।
  \item या तो ग्रेनाइट धीरे ठंडा हुआ और \omterm{def:b1:rock-forming-minerals:groups}{अभ्रक} तथा \omterm{def:b1:rock-forming-minerals:groups}{एम्फ़ीबोल} क्रिस्टलीकरण के बाद बंद हुए (शीतलन-आयुएँ), या बाद के ऊष्मन ने उनके आर्गन का एक भाग निकाल दिया। आर्गन–आर्गन आयु-स्पेक्ट्रम में एक सपाट पठार एक \omterm{def:b1:radiometric-methods:closure}{शीतलन-आयु} का समर्थन करता है; बिना पठार वाला स्पेक्ट्रम चेतावनी है कि खनिजों ने आर्गन खोया। आयुओं का क्रम, हॉर्नब्लेंड पहले और बायोटाइट बाद में, जैसा उनके \omterm{def:b1:radiometric-methods:closure}{अवरोध-तापमान} अनुमान करते हैं, भी शीतलन की ओर संकेत करता है।
  \item $\ln(1.049361)/\num{1.55125e-10} = \qty{310.6}{\Ma}$ और
$\ln(1.35784)/\num{9.8485e-10} = \qty{310.6}{\Ma}$: संगत।
  \item ज़िरकॉन \qty{310.6}{\Ma}, \qty{750}{\celsius} से ऊपर बंद: क्रिस्टलीकरण; संपूर्ण-शैल \omterm{def:b1:radiometric-methods:isochron}{समतापरेखा}, \qty[uncertainty-mode=separate]{309.6(12)}{\Ma}, भी उसे आँकती है, और अपनी अनिश्चितता के भीतर सहमत है; हॉर्नब्लेंड \qty{304.0}{\Ma}: लगभग \qty{530}{\celsius} से होकर शीतलन; बायोटाइट \qty{296.0}{\Ma}: लगभग \qty{280}{\celsius} से होकर शीतलन। ग्रेक \qty{310.6}{\Ma} पर \qty{750}{\celsius} से ऊपर से \qty{296.0}{\Ma} पर \qty{280}{\celsius} तक उतरता है।
  \item $(530 - 280)/(304.02 - 296.03) = 250/7.99 = \qty{31}{\celsius}$ प्रति दस लाख वर्ष।
  \item $(750 - 530)/(310.60 - 304.02) = 220/6.58 = \qty{33}{\celsius}$ प्रति दस लाख वर्ष कम से कम, क्योंकि ज़िरकॉन \qty{750}{\celsius} से ऊपर बंद हुआ: शीतलन आरंभ में कम से कम उतना ही तेज़ था।
  \item तापमानों का अंतर, $250 \pm \sqrt{40^2 + 40^2} = 250 \pm
        57$, \qtyrange{24}{38}{\celsius} प्रति दस लाख वर्ष देता है।
  \item $(550 - 350)/7.99 = \qty{25}{\celsius}$ प्रति दस लाख वर्ष, प्रश्न~17 की परास के भीतर: अनिश्चितता अवरोध-तापमानों से आती है, आयुओं से नहीं।
  \item $(280 - 10)/30 = \qty{9.0}{km}$।
  \item $\ln(1.039108)/\num{1.55125e-10} = \qty{247.30}{\Ma}$, और यूरेनियम-235 अनुपात से भी वही: संगत।
  \item ग्रेनाइट से, \omterm{def:b1:cross-sections:unconformity}{असंगति} पर अपरदित: उनकी आयु उसकी क्रिस्टलीकरण-आयु है। वे केवल दिखाते हैं कि बलुआ पत्थर \qty{310}{\Ma} से युवा है: एक उच्चतम आयु, राख-परत से कहीं कमज़ोर।
  \item $(58 - 18)/(247.30 - 245.30) = \qty{20.0}{m}$ प्रति दस लाख वर्ष; आधार $18/20.0 = \qty{0.90}{\Myr}$ निचली राख से पुराना है: \qty{248.20}{\Ma}।
  \item \qty{309.6}{\Ma} (\omterm{def:b1:radiometric-methods:isochron}{समतापरेखा}) या \qty{310.6}{\Ma} (ज़िरकॉन) से \qty{248.2}{\Ma} तक: लगभग \qty{62}{\Myr}।
  \item $296.0 - 248.2 = \qty{47.8}{\Myr}$ में \qty{9.0}{km}: \qty{0.19}{km} प्रति दस लाख वर्ष, यानी \qty{0.19}{mm} प्रति वर्ष। यह मानता है कि आँका गया बायोटाइट \omterm{def:b1:cross-sections:unconformity}{असंगति} के ठीक नीचे से आया, ताकि बालू के पहुँचने तक ग्रेनाइट और शिस्ट के ऊपर की सारी शैल हटा दी गई थी, और भू-तापीय प्रवणता स्थिर रही।
  \item ग्रेनाइट हॉर्नब्लेंड और बायोटाइट के बंद होने के बीच लगभग \qty{31}{\celsius} प्रति दस लाख वर्ष की दर से ठंडा हुआ।
\end{enumerate}
\end{solution}
